% COSC 3337 Final Project -- Checkpoint 1: Team Roster and Proposal
% Due Sunday, October 11, 11:59 PM. Submit the compiled PDF on Canvas.
% Length: 1--2 pages. Replace every [bracketed placeholder] and delete these comments' instructions as you go.

\documentclass[11pt]{article}
\usepackage[margin=1in]{geometry}
\usepackage{enumitem}
\usepackage{tabularx}
\usepackage{booktabs}
\usepackage{amssymb}
\usepackage{xcolor}
\usepackage[colorlinks=true,urlcolor=blue!50!black]{hyperref}
\setlist{itemsep=2pt, topsep=4pt}

% ---------- Fill these in ----------
\newcommand{\appname}{[Application Name]}
% -----------------------------------

\begin{document}

\begin{center}
  {\Large\bfseries COSC 3337 Checkpoint 1: Proposal}\\[4pt]
  {\large \appname{}}
\end{center}

\section*{1. Team Roster}
% List all 3--4 members.
\begin{enumerate}
  \item {[Full name]} --- [St.\ Edward's email]
  \item {[Full name]} --- [St.\ Edward's email]
  \item {[Full name]} --- [St.\ Edward's email]
  \item {[Full name]} --- [St.\ Edward's email]
\end{enumerate}

\section*{2. Domain and Users}
% What does the application do, who uses it, and what problem does it solve?
[One paragraph describing the application and the problem it solves.]

\medskip
\noindent\textbf{Users.} Describe each kind of user and what they do in the application.
\begin{itemize}
  \item \textbf{[User type, e.g., Buyer]:} [what they do]
  \item \textbf{[User type, e.g., Seller]:} [what they do]
\end{itemize}

\section*{3. Ten Questions}
% Plain English, from the user's point of view.
% At least 4 must combine information from more than one kind of entity (mark "Multi").
% At least 2 must involve counts, totals, or averages (mark "Agg").
\noindent\begin{tabularx}{\textwidth}{@{}rXcc@{}}
\toprule
\# & Question & Multi & Agg \\
\midrule
1  & [e.g., Which items has this seller listed that have received no bids in the past week?] & \checkmark & \\
2  & [Question] & & \\
3  & [Question] & & \\
4  & [Question] & & \\
5  & [Question] & & \\
6  & [Question] & & \\
7  & [Question] & & \\
8  & [Question] & & \\
9  & [Question] & & \\
10 & [Question] & & \\
\bottomrule
\end{tabularx}

\section*{4. Data Source}
\begin{itemize}
  \item \textbf{Source(s):} [Dataset name and URL, or ``generated with Faker/Mockaroo''; mixing is fine]
  \item \textbf{License / terms:} [License, or ``N/A -- generated'']
  \item \textbf{Expected size:} [Roughly how many records in total, and in the largest tables]
\end{itemize}

\section*{5. Entity Sketch}
% Not a formal ER diagram yet. List the main entities and how they relate.
\begin{itemize}
  \item \textbf{[Entity]:} [key attributes]
  \item \textbf{[Entity]:} [key attributes]
  \item \textbf{[Entity]:} [key attributes]
\end{itemize}
\textbf{Relationships:}
\begin{itemize}
  \item {[e.g., A seller lists many items; each item has one seller.]}
  \item {[e.g., A student enrolls in many sections; a section has many students (many-to-many).]}
\end{itemize}

\section*{6. Tech Stack}
\begin{itemize}
  \item \textbf{Database:} [e.g., MySQL on a course VM, PostgreSQL on Supabase, SQLite]
  \item \textbf{Application framework:} [e.g., Flask, Express, PHP]
  \item \textbf{Where it will run:} [Course VM / other hosting]
  \item \textbf{Git repository:} \url{[https://github.com/...]}
\end{itemize}

\end{document}
